% Compilar desde la raíz: pdflatex -output-directory=pdf_v2 temas_v2/tema1.tex
% Adaptación didáctica. Conserva los contenidos y ejercicios de temas/tema1.tex.
\input{config.tex}
% Las imágenes opcionales se sustituyen automáticamente al añadir sus archivos.
\newcommand{\imagenpendiente}[2]{%
\begin{center}\IfFileExists{imagenes/tema1/#1}{%
\includegraphics[width=.70\textwidth,height=.35\textheight,keepaspectratio]{imagenes/tema1/#1}}{%
\fbox{\parbox[c][.25\textheight][c]{.70\textwidth}{\centering\small Imagen pendiente\par\medskip #2}}}\end{center}}
\begin{document}
\begin{frame}{Funciones de varias variables}
\centering\LARGE Tema 1\par\bigskip
\large Matemáticas 2\par Grado en Ingeniería Informática
\end{frame}
\begin{frame}{Una variable}
Una función asigna una salida a cada entrada admitida.
\[f(x)=x^2\]
\[f(2)=4,\qquad f(-2)=4,\qquad f(0)=0.\]
Cada entrada tiene una única salida. Entradas distintas pueden dar la misma salida.
\end{frame}
\begin{frame}{Dos variables}
La temperatura de una habitación puede variar según la posición.
\[T=f(x,y)\]
$x$ e $y$ indican dónde medimos. $T$ indica la temperatura en ese punto.
\medskip
Fijamos la hora y la altura de la medición para estudiar solo esas dos variables.
\end{frame}
\begin{frame}{Temperatura terrestre}
A una hora determinada, la temperatura en la superficie de la Tierra depende de la longitud y la latitud.
\[T=f(x,y).\]
$x$ representa la longitud e $y$ la latitud.
\medskip
La función asigna una temperatura a cada posición admitida en el modelo.
\end{frame}
\begin{frame}{Dos entradas}
\[\underbrace{(x,y)}_{\text{entrada}}\ \longmapsto\ \underbrace{f(x,y)}_{\text{salida}}\]
Un par ordenado contiene dos datos y constituye una entrada.
\[f(x,y)=x+y,\qquad f(2,3)=2+3=5.\]
La salida sigue siendo un número real.
\end{frame}
\begin{frame}{Evaluación}
\[f(x,y)=xy\]
\[f(2,3)=6,\qquad f(0,4)=0,\qquad f(-1,3)=-3.\]
Para evaluar, sustituimos cada variable por su valor.
\medskip
El orden importa: en $g(x,y)=x+2y$, $g(1,3)=7$ y $g(3,1)=5$.
\end{frame}
\begin{frame}{Práctica breve}
Sea $f(x,y)=x+2y$.
\begin{enumerate}
\item Calcula $f(2,1)$ y $f(1,2)$.
\item Encuentra dos pares distintos cuya salida sea $4$.
\item ¿Puede un mismo par tener dos salidas?
\end{enumerate}
\end{frame}
\begin{frame}{Definición}
Una función real de dos variables reales asigna a cada par $(x,y)$ de un conjunto $D\subseteq\mathbb R^2$ un único número real.
\[f:D\longrightarrow\mathbb R,\qquad (x,y)\longmapsto f(x,y).\]
$\mathbb R^2$ es el conjunto de todos los pares de números reales.
\medskip
El conjunto $D$ se llama \textbf{dominio}.
\end{frame}
\begin{frame}{Variables}
\[z=f(x,y)\]
\begin{itemize}
\item $x$ e $y$: variables independientes, las entradas del modelo.
\item $z$: variable dependiente, la salida calculada.
\end{itemize}
En $T=f(x,y)$, la temperatura depende de la posición.
\medskip
Las entradas deben respetar las condiciones del dominio.
\end{frame}
\begin{frame}{Dominio e imagen}
\[f(x,y)=x^2+y^2\]
\textbf{Dominio}: pares que podemos introducir.
\[D=\mathbb R^2.\]
\textbf{Imagen}: valores que puede devolver la función.
\[\operatorname{Im}(f)=[0,\infty).\]
Toda salida es no negativa. Cualquier $k\ge0$ se obtiene con $(\sqrt{k},0)$.
\end{frame}
\begin{frame}{Codominio}
\[f:\mathbb R^2\longrightarrow\mathbb R,\qquad f(x,y)=x^2+y^2.\]
El \textbf{codominio} declarado es $\mathbb R$.

La \textbf{imagen} es $[0,\infty)$, pues la función nunca devuelve un valor negativo.
\medskip
Siempre se cumple $\operatorname{Im}(f)\subseteq$ codominio.
\end{frame}
\begin{frame}{Restricciones conocidas}
Trabajamos con valores reales.
\begin{center}\begin{tabular}{ll}
Expresión & Condición\\\hline
$1/x$ & $x\ne0$\\[3mm]
$\sqrt{x}$ & $x\ge0$\\[3mm]
$\ln x$ & $x>0$\\[3mm]
$\arcsin x$ & $-1\le x\le1$
\end{tabular}\end{center}
En varias variables aplicamos las mismas condiciones a la expresión completa del denominador o del argumento.
\end{frame}
\begin{frame}{Dominio sin restricciones}
\[f(x,y)=x+y\]
La suma está definida para cualquier par de números reales.
\[D=\mathbb R^2.\]
Si el enunciado no fija un dominio, buscamos el mayor conjunto donde la fórmula está definida: el \textbf{dominio natural}.
\end{frame}
\begin{frame}{Una restricción}
\[f(x,y)=\frac1{x-y}\]
El denominador debe ser distinto de cero.
\[x-y\ne0\quad\Longleftrightarrow\quad y\ne x.\]
\[D=\{(x,y)\in\mathbb R^2:y\ne x\}.\]
Admitimos todo el plano excepto la recta $y=x$.
\end{frame}
\begin{frame}{Puntos admitidos}
\[f(x,y)=\frac1{x-y}\]
\[f(2,1)=1,\qquad f(1,2)=-1.\]
El par $(2,2)$ queda fuera del dominio, porque $2-2=0$.
\medskip
Para decidir si un punto pertenece al dominio, comprobamos las restricciones antes de calcular.
\end{frame}
\begin{frame}{Dominio parabólico}
\[f(x,y)=\frac{3x-5y}{y-x^2}\]
\[y-x^2\ne0\quad\Longleftrightarrow\quad y\ne x^2.\]
El dominio es todo el plano excepto la parábola $y=x^2$.
\medskip
Un numerador nulo está permitido si el denominador no es cero.
\end{frame}
\begin{frame}{Superficie y dominio}
\[f(x,y)=\frac{3x-5y}{y-x^2}\]
\figura{s10_4.png}{.75}{.50}
La superficie representa valores sobre puntos admitidos. Sobre la parábola excluida la fórmula no define una salida.
\end{frame}
\begin{frame}{Dominio de una raíz}
\[f(x,y)=\sqrt{9-x^2-y^2}\]
El radicando debe ser no negativo.
\[9-x^2-y^2\ge0\quad\Longleftrightarrow\quad x^2+y^2\le9.\]
El dominio es el disco de centro $(0,0)$ y radio $3$, incluida su frontera.
\end{frame}
\begin{frame}{Interior y frontera}
\[f(x,y)=\sqrt{9-x^2-y^2}\]
\[f(0,0)=3,\qquad f(3,0)=0.\]
El punto $(4,0)$ queda fuera: el radicando sería $-7$.
\medskip
Sobre la frontera $x^2+y^2=9$, la salida es $0$. La imagen es $[0,3]$.
\end{frame}
\begin{frame}{Dominio del modelo}
El área de un rectángulo es $A(x,y)=xy$.
\medskip
La fórmula tiene dominio natural $\mathbb R^2$.

Si $x$ e $y$ representan longitudes de un rectángulo no degenerado, el modelo exige
\[D=\{(x,y):x>0,\ y>0\}.\]
El contexto puede restringir el dominio de una fórmula.
\end{frame}
\begin{frame}{Tres variables}
El volumen de una caja depende de largo, ancho y alto.
\[V(x,y,z)=xyz,\qquad V(2,3,4)=24.\]
Si las longitudes están en metros, el volumen está en $\mathrm{m}^3$.
\figura{s3_2.png}{.40}{.35}
\end{frame}
\begin{frame}{Puntos del espacio}
Un par $(x,y)\in\mathbb R^2$ representa un punto del plano.

Una terna $(x,y,z)\in\mathbb R^3$ representa un punto del espacio.
\medskip
Una función de tres variables asigna un número real a cada terna de su dominio.
\[f:D\subseteq\mathbb R^3\longrightarrow\mathbb R.\]
\end{frame}
\begin{frame}{Dominio logarítmico}
\[f(x,y,z)=\ln(z-y)+xy\sin z\]
El logaritmo exige $z-y>0$. El término $xy\sin z$ está definido para toda terna real.
\[D=\{(x,y,z)\in\mathbb R^3:z>y\}.\]
Es el semiespacio por encima del plano $z=y$, sin incluir el plano.
\end{frame}
\begin{frame}{Comprobación}
\[f(x,y,z)=\ln(z-y)+xy\sin z\]
$(0,1,2)$ pertenece al dominio, pues $2-1>0$.
\[f(0,1,2)=\ln1+0=0.\]
$(0,2,2)$ queda fuera, pues $\ln0$ no está definido.

$(0,3,2)$ también queda fuera, pues el argumento sería negativo.
\end{frame}
\begin{frame}{$n$ variables}
Una $n$-tupla es una lista ordenada de $n$ números:
\[(x_1,x_2,\ldots,x_n)\in\mathbb R^n.\]
Una función real de $n$ variables reales asigna a cada $n$-tupla de su dominio un único número real.
\[f:D\subseteq\mathbb R^n\longrightarrow\mathbb R,\qquad u=f(x_1,\ldots,x_n).\]
\end{frame}
\begin{frame}{Operaciones}
Para funciones $f$ y $g$ de las mismas variables:
\[(f+g)(\mathbf x)=f(\mathbf x)+g(\mathbf x),\]
\[(fg)(\mathbf x)=f(\mathbf x)g(\mathbf x).\]
Ambas operaciones requieren $\mathbf x\in D_f\cap D_g$.
\medskip
Por ejemplo, si $f(x,y)=x+y$ y $g(x,y)=xy$,
\[(f+g)(1,2)=3+2=5.\]
\end{frame}
\begin{frame}{Cociente}
\[\left(\frac fg\right)(\mathbf x)=\frac{f(\mathbf x)}{g(\mathbf x)}.\]
Además de pertenecer a ambos dominios, la entrada debe cumplir $g(\mathbf x)\ne0$.
\medskip
Con $f(x,y)=x+y$ y $g(x,y)=xy$,
\[\frac{f(x,y)}{g(x,y)}=\frac{x+y}{xy},\qquad x\ne0,\quad y\ne0.\]
\end{frame}
\begin{frame}{Composición}
La salida de una función puede ser la entrada de otra.
\[g(x,y)=x+y,\qquad h(t)=t^2.\]
\[(h\circ g)(x,y)=h(g(x,y))=(x+y)^2.\]
Para componer, $g(x,y)$ debe pertenecer al dominio de $h$.
\medskip
Si $h(t)=\sqrt t$, la composición exige $x+y\ge0$.
\end{frame}
\begin{frame}{Gráfica de una variable}
\[y=f(x)=x^2\]
Cada punto de la gráfica contiene una entrada y su salida:
\[(x,f(x)).\]
Por ejemplo, $(2,4)$ pertenece a la gráfica.
\medskip
Necesitamos dos ejes. La gráfica es una curva en el plano.
\end{frame}
\begin{frame}{Gráfica de dos variables}
\[z=f(x,y)=x+y\]
Cada punto contiene las dos entradas y la salida:
\[(x,y,f(x,y)).\]
La entrada $(1,2)$ produce el punto $(1,2,3)$.
\medskip
Necesitamos tres ejes. En este ejemplo la gráfica es un plano.
\end{frame}
\begin{frame}{El plano $z=x+y$}
\[f(0,0)=0,\qquad f(1,0)=1,\qquad f(0,1)=1.\]
\imagenpendiente{v2_plano_suma.png}{Plano $z=x+y$ con ejes $x,y,z$ y los tres puntos señalados.}
La superficie reúne todos los puntos $(x,y,x+y)$.
\end{frame}
\begin{frame}{Representación espacial}
\begin{columns}[c]
\begin{column}{.45\textwidth}
$(x,y)$ localiza la entrada en el dominio.
\medskip
$f(x,y)$ determina la altura del punto de la gráfica.
\end{column}
\begin{column}{.52\textwidth}\figura{s8_2.png}{1}{.67}\end{column}
\end{columns}
\end{frame}
\begin{frame}{Definición de gráfica}
Si $f$ tiene dominio $D\subseteq\mathbb R^2$, su gráfica es
\[G_f=\{(x,y,z)\in\mathbb R^3:(x,y)\in D,\ z=f(x,y)\}.\]
Para cada entrada admitida existe una única altura.
\medskip
Las funciones habituales de dos variables se representan mediante superficies. La forma concreta depende de la función y de su dominio.
\end{frame}
\begin{frame}{Una superficie}
\[z=f(x,y)=x+y^2,\qquad D=\mathbb R^2.\]
\figura{s9_2.png}{.85}{.52}
Con $x=0$ aparece $z=y^2$. Con $y=0$ aparece $z=x$.
\end{frame}
\begin{frame}{Una semiesfera}
\[z=\sqrt{9-x^2-y^2}.\]
\figura{s11_5.png}{.72}{.46}
Al elevar al cuadrado: $x^2+y^2+z^2=9$, con $z\ge0$.

La gráfica es la semiesfera superior. Su dominio es un disco en el plano $xy$.
\end{frame}
\begin{frame}{Más de dos entradas}
Para una función de tres variables, un punto de su gráfica contiene cuatro coordenadas:
\[(x,y,z,f(x,y,z))\in\mathbb R^4.\]
No podemos mostrar toda la gráfica con tres ejes habituales.
\medskip
Podemos estudiar cortes o representar conjuntos donde la salida tiene un valor fijo.
\end{frame}
\begin{frame}{Altura en un mapa}
\begin{columns}[c]
\begin{column}{.42\textwidth}
La altura de un terreno depende de la posición: $H=f(x,y)$.
\medskip
Una línea del mapa reúne puntos con la misma altura.
\end{column}
\begin{column}{.55\textwidth}\figura{s21_0.jpg}{1}{.67}\end{column}
\end{columns}
\end{frame}
\begin{frame}{Un nivel fijo}
Si elegimos una altura de $600\,\mathrm m$, buscamos todas las posiciones que cumplen
\[H(x,y)=600.\]
La salida queda fija. Las entradas pueden variar.
\medskip
El conjunto obtenido es un \textbf{conjunto de nivel}. En dos variables suele formar una curva.
\end{frame}
\begin{frame}{Curvas de nivel}
Para $f:D\subseteq\mathbb R^2\to\mathbb R$, el conjunto de nivel $k$ es
\[L_k=\{(x,y)\in D:f(x,y)=k\}.\]
Cuando forma una curva, hablamos de \textbf{curva de nivel}.
\medskip
Si $k\notin\operatorname{Im}(f)$, el conjunto es vacío. Algunos niveles pueden ser un solo punto.
\end{frame}
\begin{frame}{Niveles rectos}
\[f(x,y)=x+y\]
El nivel $k$ cumple $x+y=k$, es decir, $y=k-x$.
\begin{center}\begin{tabular}{ccl}
Nivel & Ecuación & Puntos del nivel\\\hline
$0$ & $y=-x$ & $(0,0)$, $(1,-1)$\\[2mm]
$2$ & $y=2-x$ & $(0,2)$, $(1,1)$\\[2mm]
$4$ & $y=4-x$ & $(0,4)$, $(2,2)$
\end{tabular}\end{center}
Las curvas de nivel son rectas paralelas.
\end{frame}
\begin{frame}{Niveles circulares}
\[f(x,y)=x^2+y^2\]
El nivel $k$ cumple $x^2+y^2=k$.
\begin{itemize}
\item $k>0$: circunferencia de radio $\sqrt k$.
\item $k=0$: únicamente el punto $(0,0)$.
\item $k<0$: conjunto vacío.
\end{itemize}
La imagen $[0,\infty)$ indica qué niveles existen.
\end{frame}
\begin{frame}{Corte horizontal}
\figura{s14_3.png}{.66}{.58}
Cortamos una superficie con el plano horizontal $z=k$.

Su proyección al plano $xy$ muestra el conjunto de nivel $f(x,y)=k$.
\end{frame}
\begin{frame}{Superficie y mapa}
\figura{s18_0.jpg}{.63}{.64}
La vista espacial muestra alturas. El mapa de niveles sitúa en el plano los puntos con igual altura.
\end{frame}
\begin{frame}{Niveles de una raíz}
\[f(x,y)=\sqrt{100-x^2-y^2}\]
El dominio es $x^2+y^2\le100$ y la imagen es $[0,10]$.
\medskip
Buscamos las posiciones cuya salida sea $c$:
\[\sqrt{100-x^2-y^2}=c.\]
Primero debe cumplirse $0\le c\le10$.
\end{frame}
\begin{frame}{Ecuación del nivel}
Para $0\le c\le10$,
\[\sqrt{100-x^2-y^2}=c\]
\[100-x^2-y^2=c^2\]
\[\boxed{x^2+y^2=100-c^2}.\]
La condición $c\ge0$ es necesaria para que elevar al cuadrado conserve la equivalencia.
\end{frame}
\begin{frame}{Radios de los niveles}
\[x^2+y^2=100-c^2\]
Para $0\le c<10$, el radio es $r=\sqrt{100-c^2}$.
\begin{center}\begin{tabular}{ccc}
Altura $c$ & Ecuación & Radio\\\hline
$0$ & $x^2+y^2=100$ & $10$\\[2mm]
$6$ & $x^2+y^2=64$ & $8$\\[2mm]
$8$ & $x^2+y^2=36$ & $6$
\end{tabular}\end{center}
Al aumentar la altura, disminuye el radio.
\end{frame}
\begin{frame}{Niveles extremos}
\[f(x,y)=\sqrt{100-x^2-y^2}\]
Para $c=10$,
\[x^2+y^2=0,\]
por lo que el nivel contiene solo $(0,0)$.
\medskip
Para $c<0$ o $c>10$, el conjunto de nivel es vacío.
\end{frame}
\begin{frame}{Vista de los niveles}
\begin{columns}[c]
\begin{column}{.45\textwidth}
\[x^2+y^2=100-c^2\]
Cada circunferencia corresponde a una altura fija.
\medskip
La cima corresponde al nivel $c=10$.
\end{column}
\begin{column}{.52\textwidth}\figura{s17_3.png}{1}{.57}\end{column}
\end{columns}
\end{frame}
\begin{frame}{Otros mapas de nivel}
\figura{s19_1.jpg}{.95}{.57}
Una superficie puede tener varios máximos. Un mismo nivel puede tener varias componentes separadas.
\end{frame}
\begin{frame}{Lectura de un mapa}
\figura{s20_1.jpg}{.95}{.52}
La superficie y sus curvas de nivel ofrecen vistas complementarias de la misma función.
\medskip
La ecuación $f(x,y)=k$ permite comprobar a qué nivel pertenece cada punto.
\end{frame}
\begin{frame}{Temperatura de un chip}
A una hora determinada, $T(x,y)$ representa la temperatura en la superficie de un procesador.
\[T(x,y)=70\]
reúne posiciones a $70\,{}^\circ\mathrm C$.
\medskip
Un mapa de niveles permite localizar zonas calientes sin representar una superficie en 3D.
\imagenpendiente{v2_chip_temperatura.png}{Mapa térmico de un procesador con escala en grados Celsius y una línea de igual temperatura.}

\end{frame}
\begin{frame}{Aplicaciones}
\begin{itemize}
\item \textbf{Medicina}: $C(x,y)$ puede representar la concentración de una sustancia en un corte de tejido. Un nivel reúne puntos con igual concentración.
\item \textbf{Meteorología}: las isobaras reúnen puntos con igual presión atmosférica.
\item \textbf{Deporte}: $t(d,v)=d/v$, con $v>0$, relaciona distancia, velocidad constante y tiempo.
\end{itemize}
En cada modelo debemos identificar entradas, salida y unidades.
\end{frame}
\begin{frame}{Superficies de nivel}
Para $f:D\subseteq\mathbb R^3\to\mathbb R$, fijamos una salida:
\[L_k=\{(x,y,z)\in D:f(x,y,z)=k\}.\]
Estos conjuntos suelen ser superficies en el espacio.
\medskip
Por ejemplo, si $f(x,y,z)=x^2+y^2+z^2$, el nivel $k=4$ es la esfera de radio $2$.
\end{frame}
\begin{frame}{Niveles cilíndricos}
\[f(x,y,z)=x^2+y^2\]
Para $k>0$,
\[x^2+y^2=k,\qquad z\in\mathbb R.\]
En cada altura aparece la misma circunferencia: el conjunto es un cilindro de radio $\sqrt k$.
\medskip
Para $k=0$ obtenemos el eje $z$. Para $k<0$, el conjunto vacío.
\end{frame}
\begin{frame}{Práctica de dominio}
Encuentra el dominio natural de cada función.
\begin{enumerate}
\item $f(x,y)=3x-y$.
\item $g(x,y)=\dfrac1{x+y}$.
\item $h(x,y)=\sqrt{4-x^2-y^2}$.
\item $p(x,y)=\ln(y-x)$.
\end{enumerate}
En cada caso, escribe la condición y describe el conjunto.
\end{frame}
\begin{frame}{Varias restricciones}
\[f(x,y)=\frac{\sqrt y}{x-1}+\ln(x+y)\]
Cada expresión aporta una condición:
\[y\ge0,\qquad x\ne1,\qquad x+y>0.\]
El dominio contiene los puntos que cumplen \textbf{todas} a la vez.
\[D=\{(x,y)\in\mathbb R^2:y\ge0,\ x\ne1,\ x+y>0\}.\]
\end{frame}
\begin{frame}{Dominio avanzado}
Encuentra el dominio de
\[f(x,y,z)=\arcsin(xy)e^{3z}+e^{(y+z)\ln(x+z)}.\]
\begin{enumerate}
\item ¿Qué condición exige el argumento de $\arcsin$?
\item ¿Qué condición exige el argumento de $\ln$?
\item ¿Añaden restricciones las exponenciales?
\end{enumerate}
Expresa el resultado como un conjunto de ternas.
\end{frame}
\begin{frame}{Comprobación del dominio}
\[f(x,y,z)=\arcsin(xy)e^{3z}+e^{(y+z)\ln(x+z)}\]
\[\boxed{D=\{(x,y,z)\in\mathbb R^3:-1\le xy\le1,\ x+z>0\}.}\]
Las exponenciales admiten cualquier exponente real definido.
\medskip
Aunque $y+z=0$, la fórmula escrita sigue exigiendo $x+z>0$ para evaluar el logaritmo.
\end{frame}
\begin{frame}{Ejercicios de nivel}
Describe los conjuntos de nivel de
\[f(x,y)=x^2+y^2,\qquad k=0,1,2,3.\]
Para cada nivel, indica la forma y, cuando corresponda, el radio.
\medskip
Después estudia
\[g(x,y,z)=x^2+y^2,\qquad k=0,1,2.\]
¿Qué cambia al permitir una tercera coordenada?
\end{frame}
\begin{frame}{Niveles hiperbólicos}
Describe los conjuntos de nivel de
\[f(x,y)=x^2-y^2,\qquad k=-2,-1,0,1,2.\]
Para $k=0$, factoriza:
\[x^2-y^2=(x-y)(x+y).\]
Para $k\ne0$, identifica las hipérbolas y su orientación.
\end{frame}
\begin{frame}{Comprobación de niveles}
Para $x^2+y^2=k$, $k=0$ da un punto y $k>0$ una circunferencia de radio $\sqrt k$.
\medskip
En tres variables, $x^2+y^2=k$ da el eje $z$ para $k=0$ y cilindros para $k>0$.
\medskip
Para $x^2-y^2=k$:
\begin{itemize}
\item $k=0$: las rectas $y=x$ e $y=-x$.
\item $k>0$: hipérbolas abiertas hacia los lados del eje $x$.
\item $k<0$: hipérbolas abiertas hacia los lados del eje $y$.
\end{itemize}
\end{frame}
\begin{frame}{Ideas clave}
\begin{itemize}
\item Una función asigna una única salida a cada entrada de su dominio.
\item El dominio reúne entradas admitidas. La imagen reúne salidas obtenidas.
\item La gráfica registra entradas y salida en un mismo punto.
\item Un conjunto de nivel reúne entradas con una salida fija.
\end{itemize}
Para estudiar una fórmula, primero identificamos sus variables y sus restricciones.
\end{frame}
\begin{frame}{Autoevaluación}
\begin{enumerate}
\item ¿Por qué $f(1,2)$ es un número, pero $(1,2)$ es una entrada?
\item ¿Por qué la gráfica de $f(x,y)$ necesita tres ejes?
\item ¿Qué diferencia hay entre dominio y nivel $k$?
\item ¿Puede un conjunto de nivel ser vacío o un punto?
\item ¿Por qué $x^2+y^2=1$ da una circunferencia en $\mathbb R^2$ y un cilindro en $\mathbb R^3$?
\end{enumerate}
\end{frame}
\end{document}
